\chapter{Conclusão}
\label{chapter:conclusion}

A duplicação do código de configuração de acessórios prejudica a manutenibilidade dos testes, mas sua remoção por meio da configuração implícita exige mais do que identificar trechos repetidos. À medida que a suíte cresce, torna-se necessário decidir como redistribuir os métodos entre classes de teste para aumentar o reuso sem introduzir mais duplicação, elevar o esforço total de desenvolvimento e manutenção ou alterar o resultado observado dos testes. Diante desse problema, este trabalho trata a distribuição dos testes como uma tarefa global de redistribuição de métodos entre classes e propõe o uso combinado de métricas de similaridade e algoritmos de agrupamento para automatizar a refatoração.

As principais contribuições propostas consistem em uma arquitetura em \textit{pipeline} e no Róża, um \textit{framework} extensível que instancia essa arquitetura com foco na configuração implícita. A separação do processo em etapas de carregamento, análise sintática, decomposição, medição, agrupamento, refatoração e escrita permite combinar implementações distintas sem acoplar todo o fluxo a uma linguagem, uma métrica ou um algoritmo específicos. O estudo sistemático apresentado neste trabalho reforça a originalidade dessa contribuição: embora existam abordagens automáticas e semiautomáticas para diferentes formas de refatoração de testes, não foi identificada uma arquitetura reusável que trate a configuração implícita como um problema global sobre toda a suíte.

No estágio atual da pesquisa, o Róża implementa as quatro primeiras etapas do \textit{pipeline} e produz a matriz de similaridade utilizada pelo agrupamento. A avaliação empírica dessa etapa comparou cinco métricas em 46 testes e 1.035 pares classificados manualmente. A LCCSS e o Simian mantiveram precisão igual a 1,0 em todos os níveis de recaptação avaliados. Entretanto, a LCCSS alcançou esse resultado sem exigir calibração e, por considerar somente o prefixo comum contíguo das configurações, representa diretamente o código que pode ser extraído para um método de configuração implícita. Esses resultados fornecem evidência favorável à etapa de medição, mas ainda não permitem confirmar a hipótese de pesquisa em sua totalidade, pois não avaliam a redistribuição global dos testes nem a redução de duplicação e de esforço após a refatoração completa.

\section{Próximos passos}

A continuação imediata do trabalho consiste em implementar no Róża as etapas de agrupamento, refatoração e escrita. Para isso, deverão ser investigados algoritmos de agrupamento e suas configurações, incluindo critérios de parada que determinem quando uma nova fusão deixa de produzir reuso suficiente para compensar a duplicação ou a fragmentação introduzida. Em seguida, a etapa de refatoração deverá reconstruir as classes, extrair os prefixos comuns para métodos de configuração implícita e preservar o comportamento observado dos testes; a etapa de escrita deverá persistir os arquivos produzidos pelo processo.

Com o \textit{pipeline} completo, será necessário avaliar a abordagem de ponta a ponta. Essa avaliação deverá medir a duplicação antes e depois da refatoração, comparar a distribuição global com alternativas restritas às classes de origem e operacionalizar o esforço de refatoração e de manutenção. Também deverá verificar se o código produzido compila e se os testes preservam seus resultados. Assim, será possível determinar em que condições a redução de duplicação compensa a reorganização das classes e, consequentemente, avaliar a hipótese de que a redistribuição global reduz a duplicação e o esforço sem alterar o comportamento observado.

A avaliação das métricas também deverá ser replicada em projetos adicionais, uma vez que os resultados atuais se restringem a testes JUnit escritos em Java e provenientes de um único sistema. A ampliação da diversidade de sistemas, domínios e níveis de teste permitirá examinar a generalização dos resultados. A extensibilidade da arquitetura também poderá ser explorada com outras métricas de similaridade e algoritmos de agrupamento.

Também se pretende estender o Róża com um algoritmo de refatoração por configuração delegada. Acredita-se que essa estratégia exigirá uma métrica de similaridade específica, capaz de identificar trechos contíguos comuns em qualquer posição da configuração de acessórios, pois, diferentemente da configuração implícita, a configuração delegada não se restringe ao prefixo compartilhado pelos testes de uma classe. A incorporação dessa métrica e desse algoritmo sem modificações nas demais etapas permitirá avaliar a flexibilidade da arquitetura para admitir diferentes níveis de fatoração do código: desde a extração de prefixos comuns para métodos de configuração implícita até a extração de trechos compartilhados por subconjuntos de testes para métodos auxiliares.